\subsection{Internal theorems for the integral construction}
\label{sec:integral-background}
\subsubsection{Decomposition and the integral domain}
\begin{theorem}[Decomposition under an envelope]\label{thm:decomposition-input}
Under the usual conditions, let $Y$ be an adapted c\`adl\`ag good integrator
with $\sup_t|Y_t|\leq q$, $q\geq0$, and $q\in L^2(Q)$.
There is a decomposition $Y-Y_0=M+A$, where $M$ is a zero-start c\`adl\`ag
local martingale and $A$ is zero-start, predictable, c\`adl\`ag and locally
of finite variation. Both components are unique up to indistinguishability. The martingale component has an exhaustive sequence of finite stopped coordinates on which it is a true $L^2$ martingale.
For bounded $S$ take a constant envelope. The good-integrator property is
preserved by equivalent measures, but the decomposition is constructed anew
under the chosen measure.
\end{theorem}
\begin{proof}[Location of the proof]
Section~\ref{sec:interface-decomposition} of the appendix derives this from
decomposition and compensation. Neither NFLVR nor integral-range closedness
is an assumption of this input.
\end{proof}

An elementary predictable test $J$ consists of finitely many stopping
intervals and bounded coefficients measurable at their left endpoints.
Its integral against any c\`adl\`ag process is a finite sum of increments.
Put
\[
 d_T^Q(Y,Z)=\sup_{J:\,|J|\leq1}
 E_Q\left[1\wedge\sup_{t\leq T}|(J\cdot(Y-Z))_t|\right].
\]
Emery convergence~\cite{Emery79} means $d_T^Q(Y_n,Y)\to0$ for every finite $T$.
A countable right-dense time set ensures measurability of the supremum.
The Cauchy quantifiers are $\forall T,\varepsilon>0\ \exists N\ \forall n,k\geq N$,
with the same $N$ for every unit-bounded test.

A \emph{local special integral} against a specially decomposed process is
obtained by applying the completed $L^2$ integral for the martingale energy
measure and the $L^1$ integral for the FV variation measure to the same
predictable coefficient, then gluing on stopped intervals. The FV integral
is the pathwise Stieltjes integral. The required localization, completion
and uniqueness are part of the next calculus theorem.

\begin{definition}[The general integral graph]\label{def:general-graph}
Write $H^{(n)}=H1_{\{|H|\leq n+1\}}$.
For predictable $H$ and zero-start strongly adapted c\`adl\`ag $X$, put
$(H,X)$ in the general integral graph of $S$ if the local special integrals
of $H^{(n)}$ converge to $X$ in Emery topology, and write $X=H\cdot S$.
Processes are identified up to indistinguishability. Neither a special
decomposition nor an infinite-time terminal value is required of $X$ itself.
\end{definition}

\begin{theorem}[Integral calculus and estimates]\label{thm:integral-input}
For bounded $S$ and the general integral graph above, the following hold.
\begin{enumerate}
\item Bounded predictable coefficients have local special integrals, unique
up to indistinguishability for a fixed coefficient. Local special integrals
belong to the general graph with the same coefficient and gain.
Every general gain is an Emery limit of elementary integrals against the
original price and is a good integrator.
\item Local special integrals obey linearity, closed stopping and predictable
restriction. For $Y=M+A$ and bounded predictable $k$,
$k\cdot Y=k\cdot M+k\cdot A$ and
$\Delta(k\cdot Y)_t=k_t\Delta Y_t$.
Finite sums use the same weights for both components. Stopping at a finite
time $T$ turns local FV into global BV.
\item Emery convergence and the general graph transfer between equivalent
probability measures with the same coefficients and gains. The same
admissibility bound and any existing intrinsic terminal value are preserved.
\item On a fixed finite horizon, a unit-bounded predictable coefficient has
unit-bounded elementary approximations simultaneously for its M and FV
integrals: terminal $L^2$ convergence on square-integrable stopped coordinates on the M side and
variation convergence on the FV side.
For a zero-start $L^2$ martingale $N$ and $|J|\leq1$,
\[
 E_Q[1\wedge\sup_{t\leq T}|(J\cdot N)_t|]\leq2\|N_T\|_2.
\]
For a local martingale $M$ exhausted by such $L^2$ stopped coordinates,
if all unit elementary tests have capped
expectation at most $b$ on that horizon, the same is true of $k\cdot M$
for every predictable $|k|\leq1$.
\item A predictable FV process $A$ has a predictable positive Hahn set $B$;
$C=1_B\cdot A$ is increasing and $C_c-C_a\geq A_c-A_a$.
For $|J|\leq1$, also
$\sup_{t\leq T}|(J\cdot A)_t|\leq\Var_{[0,T]}A$.
The intervals include their right endpoints, retaining jumps at closed stops.
\end{enumerate}
\end{theorem}
\begin{proof}[Location of the proof]
Section~\ref{sec:integral-calculus} of the appendix proves these assertions
from bounded completed integrals, truncation and uniqueness, independently
of realization of limits in the general integral domain.
\end{proof}

\subsubsection{Limit realization and original-market operations}
The next theorem supplies the M\'emin closedness result~\cite{Memin80} used in the proof of
\cite[Theorem 4.2]{DS}.
\begin{theorem}[Closedness of the general integral graph]\label{thm:graph-closure}
If general gains $Y_n$ of the same bounded price $S$ converge in Emery
topology to zero-start strongly adapted c\`adl\`ag $X$, some predictable $H$
satisfies $(H,X)$ in the same general integral graph.
\end{theorem}
\begin{proof}[Location of the proof]
Section~\ref{sec:closure-proof} of the appendix connects topology comparison,
common localization and realization of both components by one coefficient.
The subsequent argument uses only this conclusion, not the auxiliary
measure or stopping sequence selected in its proof.
\end{proof}

\begin{corollary}[Operations in the original market]\label{cor:market-operations}
Finite linear combinations, closed stops, restrictions to stopping intervals,
and finite predictable pastings of general gains are again general gains
of the original price. Moreover, if a realized zero-start gain $Y=M+A$
has an $L^2(Q)$ envelope, every realized local special integral $V=k\cdot Y$
is a general gain of the original price. Its components remain the same
$V=k\cdot M+k\cdot A$. If $Q$ is equivalent to the original measure, the
same gain can be returned to that measure. If $V\geq-a$ and $V$ is constant
after $T$, its intrinsic terminal value is $V_T$, in the terminal set at
the same admissibility level.
\end{corollary}
\begin{proof}
Apply finite sums, stops and restrictions to the elementary approximations
of Theorem~\ref{thm:integral-input}. Finite-sum associativity and the triangle
inequality give Emery convergence to the specified transformed gain.
Pasting coefficients multiply events known at a stopping time by indicators
of subsequent stopping intervals, so are predictable.
Theorem~\ref{thm:graph-closure} realizes these regular zero-start limits.

For bounded $k$, approximate the M and FV integrals with the same elementary
coefficients. Each elementary transform is a finite sum of stopped increments
of $Y$; substitute its original-price approximations. First make the error
small in each finite sum, then diagonalize over integer horizons. This gives
original-price elementary integrals converging to the same $V$ in Emery
topology. Apply closedness again. For unbounded realized $k$, truncation
convergence from Theorem~\ref{thm:integral-input} reduces to the bounded case;
apply the same closedness theorem. Transfer the general graph between measures.
The bound, constancy after $T$, and terminal value are properties of the
unchanged gain and are preserved.
\end{proof}
